    We then explore an extreme example of inference issues. \cite{chen2018effect} used an innovative fuzzy regression discontinuity design to estimate the effects of air quality alerts on health outcomes in Toronto, Canada. Their units are 143 daily observations at the city-level over the 2003–2012 period. The authors found that "alert announcements reduced asthma-related emergency department visits by 4.73 $\pm$ 2.37 cases per 1 000 000 people per day". It represents in relative terms a large decrease of 25\% in emergency department visits for the pathology. According to us, it represents a large decrease which authors do explain with their available data---is it due to a drop in air pollution or people avoidance behavior?---and do not put into perspective with other air pollution interventions or seasonal variation in emergency department visits. For instance, it would interesting to see how asthma-related emergency department visits evolve between week days and weekend or on bank days and holidays.
        
        \begin{figure}[h!]
            \caption{Power, Type M and S Errors Curves for \cite{chen2018effect}}
            \label{fig:chen}
        \centering\includegraphics[width=\linewidth]{2.literature_review/graph_chen_iv_2018.pdf}
            \caption*{\footnotesize \textit{Notes}: In each panel, a metrics, such as the statistical power, the exaggeration ratio or the probability to make a type S error, is plotted against a range of hypothetical effect sizes. The solid orange line represents the observed two-stage least square estimate reported in the article.}
        \end{figure}
        
    In \Cref{fig:chen}, we plot the power, type M and S errors curves of the study for a range of hypothetical effect sizes. The solid orange line represents the observed two-stage least square estimate reported in the article. Imagine that the true decrease in daily emergency-department visits for asthma was 2.4 - half the estimate they found (the dashed orange line in \Cref{fig:chen}). The magnitude of the effect would be still substantial as it would represent a 12.5\% increase in asthma visits. The statistical power would be 17\% and the overestimation factor would be equal to 2.5. For this study, the type M error is very worrying. The study highlights a particular tension in the causal inference literature where researchers sometimes work with more credible methods which exploit a fraction of the data at the cost of decreasing their statistical power and running into type M errors.
    
    
     Articles based on causal inference methods rarely compute uncertainty intervals and discuss the range of effect size supported by the data---68\% of the article explicitly dichotomize evidence between statistically significant and not statistically significant estimates. Nearly all published estimates are statistically significant. We find some clues about the possibility of publication bias since most \textit{t}-statistics are gathered nearby 1.96, the threshold for an estimate to be statistically significant at the 5\% level. In only 8\% of the articles the term "statistical power" is mentioned by researchers, often as a way to justify their non-statistically significant estimates. 
     
     
     
    %Researchers usually proceed as follows: they first define the causal estimand they are interested in, then run a model to retrieve the estimate and standard error, implement several sensitivity analyses and placebo tests, and finally interpret the results. These are the standard procedures but more could be done by researchers to gauge the credibility of their findings. Ideally, they should replicate several times their research design to evaluate how estimates would vary. It would be however impossible to ask such a task as researchers investigating the acute effects of air pollution on health rely on administrative data. They often have little margins to control the spatial and temporal scopes of the data set. Fortunately, it is possible to carry out hypothetical replication of their research design using only the range of effect sizes that they would deem plausible for the treatment of interest and the standard error of their estimates.
